\chapter{Artefakt 1 - Automatisierung von Logistics Equipment Assemblies}
\label{chap:Art1LEA}

Dieses Kapitel stellt die \emph{MTP-basierte Automatisierung und Integration von Logistics Equipment Assemblies} als erstes Artefakt dieser Dissertation vor. Dabei werden sowohl bereits etablierte als auch neu spezifizierte MTP-Konzepte behandelt. Teile dieser Konzepte wurden bereits in \cite{Blumenstein.atpLEA,Blumenstein.MoProLog} veröffentlicht; ergänzend untersuchte die studentische Arbeit \cite{Janzen.2023} konkret Parametrierungsmechanismen für LEAs. Die Evaluationsbeispiele in Abschnitt~\ref{sec:Evaluationsbeispiele} veranschaulichen die Anwendung des Artefakts auf industrienahe Anwendungsfälle der produktionsnahen Logistik. Darüber hinaus enthält das GitHub-Repository zu dieser Dissertation \cite{Blumenstein.Github} Detailspezifikationen der eingeführten MTP-Erweiterungen sowie eine noch detailliertere Beschreibung ihrer Anwendung auf die Automatisierung einer Palettier-LEA und einer Stretchhauben-LEA.

\section{Artefaktübersicht}
\label{sec:Art1LEAUebersicht}

Die MTP-basierte Automatisierung von LEAs besteht aus den gleichen Bestandteilen wie die Automatisierung von PEAs. Diese sind in Abbildung~\ref{fig:UebersichtLEA} dargestellt und in den nachfolgenden Abschnitten detailliert beschrieben. 

\begin{figure}[htbp]
	\centerline{\includegraphics[width=0.4\textwidth]{Inhalt/Abbildungen/06_LEA-Automatisierung/Uebersicht.pdf}}
	\caption{Bestandteile der LEA-Automatisierung}
	\label{fig:UebersichtLEA}
\end{figure}

Die Grundlage bildet eine \textbf{dienstebasierte Automatisierung} mit \underline{einem} MTP-Dienst pro LEA (Abschnitt~\ref{sec:LeaDienste}). Dieser kann als auftragsorientierter Cyclic Execution Service oder bedarfsorientierter Single Execution Service ausgeführt werden, um die in Abschnitt~\ref{subsec:ArbeitsweiseMLS} beschriebenen LEA-Arbeitsweisen umzusetzen. Für die \textbf{Parametrierung} dieser Dienste (Abschnitt~\ref{sec:LeaParameter}) stehen je nach Art der Parameter verschiedene Parametrierungsmechanismen bereit. Für auftragsspezifische und konstruktionsspezifische Parameter werden heutige Mechanismen der MTP-Spezifikation genutzt. Für produkt- und verpackungsspezifische Parameter wird eine Parameterverwaltung in LEAs eingeführt, die eine Verwaltung und Übertragung von Parametersätzen ermöglicht. Über die bestehende MTP-Spezifikation hinaus unterstützt dieses Artefakt strukturierte und arraybasierte Parameter. Für die \textbf{Reportwerte} (Abschnitt~\ref{sec:Reportwerte}) und \textbf{Prozesswerte} (Abschnitt~\ref{sec:Prozesswerte}) eines LEA-Dienstes werden die in \cite{PNO.2025.Part4} spezifizierten Konzepte genutzt und diese ähnlich wie bei den Parametern um strukturierte und arraybasierte Datentypen erweitert. Die Visualisierung der Prozesse in einer LEA erfolgt mittels \textbf{LEA-Bedienbildern} (Abschnitt~\ref{sec:LeaHmi}). Neben den in der MTP-Spezifikation bereits vorgesehenen Bedienbildkonzepten wird die Möglichkeit eingeführt, LEA-spezifische statische Bedienbildobjekte einzufügen und Variablen mit strukturierten und arraybasierten Datentypen einzubinden.

\section{Dienstebasierte Automatisierung}
\label{sec:LeaDienste}

Im Kontext der dienstebasierten LEA-Automatisierung ist es sinnvoll, die Funktionalität einer jeden LEA (im Gegensatz zu PEAs) in nur einem einzigen (Haupt-) Dienst im Sinne des MTP-Konzepts umzusetzen. Dies liegt darin begründet, dass LEAs einen vordefinierten physischen Aufbau mit Achsen besitzen, die verfahren werden können. Die Geschwindigkeit und die Reihenfolge der Ansteuerung dieser Achsen kann (durch Parametrierung) variiert werden. Allerdings können keine völlig verschiedenen Bewegungsmuster und damit verschiedenen Funktionen durch eine LEA umgesetzt werden. Somit stellt jede LEA eine bestimmte Hauptfunktionalität bereit, die vor allem durch Parametrierung angepasst wird \cite{Blumenstein.Designprinzipien}. Dies bestätigt auch die betrachtete Literatur, v.~a. \cite{InternationalSocietyofAutomation.2022,GrunerMTPPackMLatp.2022}. Eine interne Dienstehierarchie unterhalb des (Haupt-) Dienstes ist, z.~B. durch die Anbindung von FEAs, möglich, wird aber an dieser Stelle nicht weiterführend betrachtet. 

Bei der Umsetzung der LEA-Dienste besteht gemäß REQ-LEA-1b die Anforderung nach einer auftrags- und bedarfsorientierten Arbeitsweise. Für diese beiden Arbeitsweisen werden im Folgenden zwei Arten der Diensteausführung für LEA-Dienste eingeführt -- der auftragsorientiert arbeitende \textbf{\gls{ces}} und der bedarfsorientiert arbeitende \textbf{\gls{ses}}. Diese werden in den Abschnitten~\ref{sec:CES} und \ref{sec:SES} detailliert beschrieben.

Für einige LEA-Typen ist es sinnvoll, ihre Funktionalität sowohl im CES- als auch im SES-Betrieb anzubieten. So können sie in verschiedenen MLSs wahlweise im CES- oder SES-Betrieb betrieben werden, aber nicht beides gleichzeitig. Aufgrund dessen werden CES- und SES-Betrieb als verschiedene Prozeduren eines LEA-Dienstes umgesetzt. Eine LEA kann zudem für den Einsatz in verschiedenen MLS-Kontexten auch mehrere CES- und SES-Prozeduren aufweisen. CES- und SES-Prozeduren sind Prozeduren konform zur \cite{PNO.2025.Part4}, allerdings liegen ihnen spezielle Interpretationen des MTP-Zustandsautomaten zugrunde, die in den Abschnitten~\ref{sec:CES} und \ref{sec:SES} beschrieben werden. 

Um CES- und SES-Prozeduren für einen LOL unterscheidbar zu machen, ist für jede Prozedur eines LEA-Dienstes maschinenlesbar zu hinterlegen, um welche der beiden Arten der Diensteausführung es sich handelt. Zu diesem Zweck werden \emph{FunctionClassificationAttributes} gemäß \cite{PNO.2025.Part4} vorgesehen. 

\subsection{Cyclic Execution Service}
\label{sec:CES}

Der Cyclic Execution Service (CES) zur auftragsorientierten Bearbeitung von LOs ist dafür konzipiert, einen Auftrag entgegenzunehmen und anschließend alle LOs dieses Auftrags in gleicher Weise zu bearbeiten. Charakteristisch für einen Dienst im CES-Betrieb ist, dass dieser einmal zu Beginn seiner Ausführung gemäß den Auftragsdaten parametriert wird und dann zyklisch eine bestimmte oder unbestimmte Anzahl gleichartiger LOs bearbeitet. Dies ermöglicht eine schnelle Bearbeitung und einen geringen Kommunikationsaufwand zum übergeordneten LOL.

Abbildung~\ref{fig:CesZa} zeigt die Interpretation der Hauptschleife des MTP-Zu\-stands\-auto\-ma\-ten für einen Dienst im CES-Betrieb. Der Zustand einer CES-Prozedur re\-prä\-sen\-tiert den Zustand der LEA und nicht den Zustand der Bearbeitung eines LO (Abbildung~\ref{fig:CesZa}, orange Kästchen). Die Bearbeitung der LOs läuft innerhalb des EXE\-CU\-TE-Zu\-stands ab (Abbildung~\ref{fig:CesZa}, grüne Pfeile). 

\begin{figure}[htbp]
	\centerline{\includegraphics[width=0.9\textwidth]{Inhalt/Abbildungen/06_LEA-Automatisierung/Dienste/CES_ZA.pdf}}
	\caption[Interpretation des MTP-Zustandsautomaten für einen Dienst im \newline CES-Betrieb]{Interpretation des MTP-Zustandsautomaten für einen Dienst im CES-Betrieb}
	\label{fig:CesZa}
\end{figure}

Abbildung~\ref{fig:CesInteraktion} zeigt die Interaktionen eines Dienstes im CES-Betrieb mit einem übergeordneten LOL während eines Dienstedurchlaufs. Wie jeder MTP-Dienst befindet sich ein LEA-Dienst im CES-Betrieb zu Beginn seiner Ausführung im Initialzustand IDLE. In diesem Zustand werden dem Dienst alle für die Ausführung notwendigen Auftragsdaten durch ein Auftrags-Management im LOL oder durch einen Operator übergeben. Während der späteren Ausführung ist im normalen Betrieb keine Umparametrierung mehr vorgesehen. Zur Übergabe der Auftragsdaten wird eine Parametrierung gemäß Abschnitt~\ref{sec:LeaParameter} vorgenommen. 

\begin{figure}[H]
	\centerline{\includegraphics[width=0.5\textwidth]{Inhalt/Abbildungen/06_LEA-Automatisierung/Dienste/Arbeitsweise_CES.pdf}}
	\caption{Interaktion eines Dienstes im CES-Betrieb mit einem übergeordneten Logistics Orchestration Layer}
	\label{fig:CesInteraktion}
\end{figure}

Nachdem die CES-Prozedur ihre Startbereitschaft signalisiert hat (StartEn = true), kann ein \emph{Start}-Befehl abgesetzt werden und die LEA fährt im Zustand STARTING hoch. Im nachfolgenden EXECUTE-Zustand erfolgt eine zyklische Bearbeitung von LOs auf Basis der zuvor eingestellten Auftragsdaten. So werden alle LOs des Auftrags auf die gleiche Weise verpackt. Je nachdem, ob eine kontinuierliche oder eine selbstbeendende CES-Prozedur vorliegt, wird die Bearbeitung durch einen \emph{Complete}-Befehl oder bei Eintreten einer bestimmten Bedingung (z.~B. definierte Anzahl bearbeiteter LOs) beendet. In COMPLETING werden anschließend alle noch in der Bearbeitung befindlichen LOs fertiggestellt und die LEA leergefahren. Durch den Zustand COMPLETED wird die vollständige Beendigung der Bearbeitung signalisiert. Durch einen \emph{Reset}-Befehl wird die Prozedur schließlich wieder in den Zustand IDLE überführt und kann einen neuen Auftrag entgegennehmen. Die Pause-, Hold-, Stop- und Abort-Schleife können entsprechend den in \cite{PNO.2025.Part4} beschriebenen Konventionen durchlaufen werden. 

\subsection{Single Execution Service}
\label{sec:SES}

Der Single Execution Service (SES) zur bedarfsorientierten Bearbeitung von LOs ist dafür konzipiert, LOs auf verschiedene Art und Weise gemäß ihrer individuellen Auftragsdaten zu bearbeiten. Charakteristisch für einen Dienst im SES-Betrieb ist, dass dessen Parametrierung für jedes zu bearbeitende LO spezifisch angepasst wird. Dies ermöglicht eine flexible Bearbeitung von LOs, die zu unterschiedlichen Aufträgen gehören und/oder einen unterschiedlichen Bearbeitungszustand aufweisen.
 
Abbildung~\ref{fig:SesZa} zeigt die Interpretation der Haupt- und Pauseschleife des MTP-Zustandsautomaten für einen Dienst im SES-Betrieb. Die orange-markierten Zu\-stän\-de bilden ähnlich wie bei CES-Prozeduren den aktuellen Zustand der LEA ab. Die grün markierten Zustände hingegen bilden den Zustand der Bearbeitung eines ggf. in der LEA befindlichen LOs ab.

\begin{figure}[htbp]
	\centerline{\includegraphics[width=0.9\textwidth]{Inhalt/Abbildungen/06_LEA-Automatisierung/Dienste/SES_ZA.pdf}}
	\caption[Interpretation des MTP-Zustandsautomaten für einen Dienst im \newline SES-Betrieb]{Interpretation des MTP-Zustandsautomaten für einen Dienst im SES-Betrieb}
	\label{fig:SesZa}
\end{figure}

Abbildung~\ref{fig:SesInteraktion} zeigt die Interaktionen eines Dienstes im SES-Betrieb mit einem übergeordneten LOL während eines Dienstedurchlaufs. 

\begin{figure}[H]
	\centerline{\includegraphics[width=0.65\textwidth]{Inhalt/Abbildungen/06_LEA-Automatisierung/Dienste/Arbeitsweise_SES.pdf}}
	\caption{Interaktion eines Dienstes im SES-Betrieb mit einem übergeordneten Logistics Orchestration Layer}
	\label{fig:SesInteraktion}
\end{figure}
 
Auch ein LEA-Dienst im SES-Betrieb befindet sich zu Beginn seiner Ausführung im Initialzustand IDLE. In diesem Zustand kann dem Dienst ein Teil der notwendigen Parameter mittels der Parametrierungsmechanismen aus Abschnitt~\ref{sec:LeaParameter} übergeben werden. Zu diesem Zeitpunkt steht allerdings noch nicht fest, welche LOs in welcher Reihenfolge beim SES eintreffen. Daher können nur Parameter übergeben werden, die unabhängig von der Art und dem Bearbeitungszustand des zu bearbeitenden LO sind, z.~B. den konstruktiven Aufbau der LEA betreffende Konfigurationsparameter. Zudem können bereits Parametersätze für verschiedene LO-Arten und -Bearbeitungszustände an die SES-Prozedur übergeben werden, die im Folgenden ausgewählt werden, wenn ein entsprechendes LO eintrifft. Nachdem die SES-Prozedur ihre Startbereitschaft signalisiert hat (StartEn = true), kann ein \emph{Start}-Befehl abgesetzt werden und die LEA fährt im Zustand STARTING ohne einen konkreten Auftrag hoch. Anschließend wechselt die SES-Prozedur über die Zustände EXECUTE und PAUSING in den Zustand PAUSED. Dieser Zustand signalisiert, dass der SES aktiv ist und auf einen externen Trigger wartet, der ihm signalisiert, dass ein LO bearbeitet werden soll. Dies kann z.~B. die Ankunft eines AGV sein, welches den Auftrag hat, ein LO an die betreffende LEA zu übergeben. Bei Erkennung eines solchen Triggers wechselt der Zustand des LEA-Dienstes in RESUMING und die Art (\emph{ProductId}) sowie der Bearbeitungszustand (\emph{LogisticsObjectStatus}) des zu bearbeitenden LO werden identifiziert. Daraufhin erfolgt die Parametrierung der SES-Prozedur gemäß den individuellen Auftragsdaten des LO. Im nachfolgenden EXECUTE-Zustand wird die Bearbeitung des einzelnen LO bedarfsgerecht ausgeführt. Nach erfolgter Bearbeitung wechselt der SES wieder über den PAUSING- in den PAUSED-Zustand und wartet auf den nächsten Trigger zur Bearbeitung eines LO. SES-Prozeduren werden immer kontinuierlich ausgeführt, da zu Beginn der Diensteausführung nicht feststeht, wie viele LOs in welcher Reihenfolge bearbeitet werden müssen. Sind keine weiteren LOs mehr zu bearbeiten, ist die SES-Prozedur mittels eines \emph{Complete}-Befehls zu beenden. Ist zu diesem Zeitpunkt noch ein LO in der LEA, wird dessen Bearbeitung im Zustand COMPLETING abgeschlossen und die LEA leergefahren. Durch den Zustand COMPLETED wird die vollständige Beendigung der Bearbeitung signalisiert. Durch einen \emph{Reset}-Befehl wird die Prozedur schließlich wieder in den Zustand IDLE überführt und die LEA kann gestartet werden, wenn wieder Bedarf besteht. Die Hold-, Stop- und Abort-Schleife können entsprechend den in \cite{PNO.2025.Part4} beschriebenen Konventionen durchlaufen werden. 

\section{Parametrierung}
\label{sec:LeaParameter}

In diesem Abschnitt werden MTP-basierte Konzepte zur Parametrierung von LEA-Diensten vorgestellt. In Abschnitt~\ref{subsec:Parameterarten} werden zunächst zu übertragende Parameterarten untersucht. Die nachfolgenden Abschnitte~\ref{subsec:ParameterUebertragung} und \ref{subsec:ParameterInitiierung} stellen Grundbausteine der LEA-Parametrierung vor. In Abschnitt~\ref{subsec:Parametrierngsmech} wird darauf basierend ein Parametrierungsmechanismus für die verschiedenen Parameterarten von LEA-Diensten empfohlen. 

\subsection{Arten von LEA-Parametern}
\label{subsec:Parameterarten}

Im Kontext der Parametrierung von LEA-Diensten lassen sich die in Tabelle~\ref{tab:Parameterarten} dargestellten Parameterarten unterscheiden \cite{Blumenstein.atpLEA,Janzen.2023,Blumenstein.MoProLog}. 

% Parameterarten
\begin{footnotesize}
	\begin{longtable}{|
		>{\columncolor[HTML]{E0E0E0}}l 
		>{\columncolor[HTML]{FFFFFF}}l |}
		\caption{Parameterarten von Logistics Equipment Assemblies}
		\label{tab:Parameterarten}\\
		\hline

		\multicolumn{1}{|L{2.5cm}|}{\cellcolor[HTML]{E0E0E0} \textbf{Bezeichnung}}        	& 
		\multicolumn{1}{L{12.13cm}|}{\cellcolor[HTML]{E0E0E0}\textbf{Beschreibung}} 
		\\ \hline
	
		\multicolumn{1}{|L{2.5cm}|}{\cellcolor[HTML]{FFFFFF} Auftrags-spezifische\newline Parameter }        	& 
		\multicolumn{1}{L{12.13cm}|}{\cellcolor[HTML]{FFFFFF} Diese Parameter dienen der Übergabe von Auftragsdaten an den LEA-Dienst. Sie resultieren aus Kundenaufträgen und ändern sich mit jedem Auftrag. Sie geben die organisatorischen Daten des Auftrags (z.~B. die Auftragsnummer), das zu fertigende Produkt und dessen Menge (z.~B. die Sack- oder Palettenanzahl) vor. Parameter dieser Art sind ausschließlich bei auftragsorientierter Arbeitsweise, d.~h. bei Prozeduren im CES-Betrieb, vorhanden.} 
		\\ \hline

		\multicolumn{1}{|L{2.5cm}|}{\cellcolor[HTML]{FFFFFF} Produkt-spezifische\newline Parameter }        	& 
		\multicolumn{1}{L{12.13cm}|}{\cellcolor[HTML]{FFFFFF} Diese Parameter ergeben sich aus dem zu verpackenden LO inkl. seiner kunden- und länderspezifischen Besonderheiten (im Folgenden \emph{Produkt} genannt). Sie müssen angepasst werden, wenn ein anderes Produkt gefertigt werden soll. Beispiele sind Stretchparameter oder Packmuster.} 
		\\ \hline

		\multicolumn{1}{|L{2.5cm}|}{\cellcolor[HTML]{FFFFFF} Verpackungs-spezifische\newline Parameter }        	& 
		\multicolumn{1}{L{12.13cm}|}{\cellcolor[HTML]{FFFFFF} Diese Parameter ergeben sich aus dem zu verwendenden Verpackungsmaterial (z.~B. Paletten- oder Folientyp). Sie müssen angepasst werden, wenn ein anderes Verpackungsmaterial für ein Produkt genutzt werden soll. Beispiele sind die Abmaße einer Palette oder die Dicke und das Material einer Folie.} 
		\\ \hline

		\multicolumn{1}{|L{2.5cm}|}{\cellcolor[HTML]{FFFFFF} Konstruktions-spezifische\newline Parameter }        	& 
		\multicolumn{1}{L{12.13cm}|}{\cellcolor[HTML]{FFFFFF} Diese Parameter sind bedingt durch den physischen Aufbau der LEA oder bilden grundlegende Einstellungen (z.~B. Geschwindigkeitsvorgabe) unabhängig vom Verpackungsauftrag ab. Sie ändern sich z.~B. beim Umbau bzw. der Bestückung der LEA. Sie können z.~B. vorgeben, welche FEAs einer LEA untergeordnet sind, oder mit welchen Verbrauchsstoffen (z.~B. Paletten) die LEA bestückt ist.} 
		\\ \hline

		
	\end{longtable}
\end{footnotesize}

\subsection{Übertragung von LEA-Parametern}
\label{subsec:ParameterUebertragung}

Im Kontext der Parametrierung von LEA-Diensten lassen sich die in Tabelle~\ref{tab:ParameterUebertragung} dargestellten Mechanismen zur Parameterübertragung unterscheiden \cite{Blumenstein.atpLEA,Janzen.2023,Blumenstein.MoProLog}. Nachfolgend wird beschrieben, welche Vor- und Nachteile diese Mechanismen aufweisen und wie sie mit bestehenden und neuen MTP-Konzepten umgesetzt werden können.

% Bewertung Parametrierungsmechanismen
\begin{footnotesize}
	\begin{longtable}{|
		>{\columncolor[HTML]{E0E0E0}}l 
		>{\columncolor[HTML]{FFFFFF}}l |}
		\caption{Arten der Übertragung von Parametern an Logistics Equipment Assemblies}
		\label{tab:ParameterUebertragung}\\
		\hline

		\multicolumn{2}{|L{14.8cm}|}{\cellcolor[HTML]{E0E0E0}\textbf{Übertragung von Einzelparametern}} 
		\\ \hline

		\multicolumn{2}{|L{14.8cm}|}{\cellcolor[HTML]{FFFFFF}{Dieser Mechanismus sieht vor, dass die LEA-Parameter als Einzelwerte mit primitiven Datentypen über separate Parameterschnittstellen an den LEA-Dienst übertragen werden.}} 
		\\ \hline

		\multicolumn{1}{|L{1.5cm}|}{\cellcolor[HTML]{FFFFFF} Vorteile }        	& 
		\multicolumn{1}{L{13.13cm}|}{\cellcolor[HTML]{FFFFFF} Jedem Parameter können Metainformationen (z.~B. Minimal-/Maximalwert, Einheit) mitgegeben werden. Nur die aktuell verwendeten Parameterwerte müssen abgespeichert werden, somit ist der notwendige Speicherbedarf gering.} 
		\\ \hline

		\multicolumn{1}{|L{1.5cm}|}{\cellcolor[HTML]{FFFFFF} Nachteile }        	& 
		\multicolumn{1}{L{13.13cm}|}{\cellcolor[HTML]{FFFFFF} Zur Parametrierung von LEAs können sehr viele Parameter notwendig sein, wodurch die Diensteschnittstelle umfangreich und die Parametrierung aufwendig wird. Zudem wird jeder Parameter einzeln an den LEA-Dienst übergeben. Somit muss stets sichergestellt werden, dass ein konsistenter, valider Datensatz über alle Parameter hinweg am Dienst anliegt. Bei einem Ausfall des Parameter-Managements ist keine automatisierte Parametrierung mehr möglich.} 
		\\ \hline

		\multicolumn{2}{|L{14.8cm}|}{\cellcolor[HTML]{E0E0E0}\textbf{Übertragung einzelner Parametersätze}} 
		\\ \hline

		\multicolumn{2}{|L{14.8cm}|}{\cellcolor[HTML]{FFFFFF}{Dieser Mechanismus sieht vor, dass die LEA-Parameter als Parametersatz mit einem LEA-spezifischen strukturierten Datentyp über eine oder wenige Parameterschnittstellen an den LEA-Dienst übertragen werden.}} 
		\\ \hline

		\multicolumn{1}{|L{1.5cm}|}{\cellcolor[HTML]{FFFFFF} Vorteile }        	& 
		\multicolumn{1}{L{13.13cm}|}{\cellcolor[HTML]{FFFFFF} Vor allem bei LEAs mit großen Parametersätzen wird die Diensteschnittstelle vereinfacht und der Aufwand zur Parametrierung verringert. Zudem sind ein konsistentes Schreiben und Übernehmen des kompletten Parametersatzes möglich. Nur die aktuell verwendeten Parameterwerte müssen abgespeichert werden, somit ist der notwendige Speicherbedarf gering.} 
		\\ \hline

		% zu den Metadaten: Innerhalb einer Struktur können nicht Variablen lesend und schreibens sein, während andere nur lesend sind. Wenn ich die ganze Struktur schreibe, würde ich dann auf Variablen schreiben, die nur lesend sind. OPC UA hat Mechanismen drin, das zu ermöglichen, aber das ist technisch unsauber.

		\multicolumn{1}{|L{1.5cm}|}{\cellcolor[HTML]{FFFFFF} Nachteile }        	& 
		\multicolumn{1}{L{13.13cm}|}{\cellcolor[HTML]{FFFFFF} Den einzelnen Parametern des Parametersatzes können keine Metainformationen mitgegeben werden. Die Parameter würden so eine Mischung aus rein lesbaren sowie les"~/schreibbaren Variablen enthalten. Beim Schreiben auf ein solches strukturiertes Element würde demnach auch auf rein lesbare Variablen geschrieben werden, was technisch unsauber ist. Darüber hinaus muss bei der Übertragung von Parametersätzen stets der gesamte Parametersatz anstatt einzelner Variablen übertragen werden, was zu einer hohen Netzwerklast führen kann. Bei einem Ausfall des Parameter-Managements ist keine automatisierte Parametrierung mehr möglich.} 
		\\ \hline

		\multicolumn{2}{|L{14.8cm}|}{\cellcolor[HTML]{E0E0E0}\textbf{Übertragung einer Auswahl an Parametersätzen}} 
		\\ \hline

		\multicolumn{2}{|L{14.8cm}|}{\cellcolor[HTML]{FFFFFF}{Dieser Mechanismus sieht vor, dass Parametersätze für verschiedene Produkte an eine LEA übertragen werden und in einer LEA-internen Datenhaltung verwaltet werden. Über eine ID kann anschließend je nach zu verpackendem Produkt ausgewählt werden, welcher Parametersatz für den aktuellen Verpackungsvorgang verwendet werden soll. Die Parametersätze weisen einen LEA-spezifischen strukturierten Datentyp auf.}} 
		\\ \hline

		\multicolumn{1}{|L{1.5cm}|}{\cellcolor[HTML]{FFFFFF} Vorteile }        	& 
		\multicolumn{1}{L{13.13cm}|}{\cellcolor[HTML]{FFFFFF} Es ist eine schnelle und einfache Auswahl des zu nutzenden Parametersatzes mit sehr geringem Kommunikationsaufwand möglich. Zudem ist stets die Konsistenz der Parametersätze sichergestellt. Durch das Vorhalten von Parametersätzen ist eine automatisierte Parametrierung unter Umständen auch bei einem Ausfall des Parameter-Managements möglich.} 
		\\ \hline

		\multicolumn{1}{|L{1.5cm}|}{\cellcolor[HTML]{FFFFFF} Nachteile }        	& 
		\multicolumn{1}{L{13.13cm}|}{\cellcolor[HTML]{FFFFFF} Es sind zwei Schnittstellen notwendig – eine zum Laden der Parametersätze in die LEA und eine für die ID-basierte Auswahl eines Parametersatzes. Es besteht noch keine Möglichkeit, den Zusammenhang zwischen beiden Schnittstellen zu modellieren. Da verschiedene Parametersätze vorgehalten werden, ist der notwendige Speicherbedarf höher als bei den anderen beiden Mechanismen.} 
		\\ \hline
		
	\end{longtable}
\end{footnotesize}

Diese Arten der Parameterübertragung sind gemäß REQ-LEA-1e und REQ-\linebreak MLS-1 im Kontext der LEA-Automatisierung MTP-basiert umzusetzen.

\textbf{Übertragung von Einzelparametern:} Dieser Mechanismus entspricht der Parametrierung, die in den bisherigen MTP-Konzepten vorgesehen ist. Somit stehen bereits entsprechende Parameterschnittstellen für die Datentypen REAL, DINT, BOOL und STRING in \cite{PNO.2025.Part4} zur Verfügung. Es ist keine Erweiterung der MTP-Spezifikation notwendig.

\textbf{Übertragung einzelner Parametersätze:} Bisher sind im MTP-Konzept keine Parameterschnittstellen für strukturierte Datentypen vorgesehen. Daher wird eine \emph{StructServParam}-Schnittstelle eingeführt, die die Übergabe eines Parametersatzes mit LEA-spezifischem strukturiertem Datentyp an den LEA-Dienst ermöglicht. Zudem wird eine Möglichkeit beschrieben, die notwendigen komplexen Datentypen im MTP sowie im OPC~UA-Server einer LEA zu modellieren.

\textbf{Übertragung einer Auswahl an Parametersätzen:} Zur MTP-basierten Umsetzung der Schnittstelle für die ID-Auswahl kann auf die in \cite{PNO.2025.Part4} spezifizierte \emph{DIntServParam}-Schnittstelle zurückgegriffen werden. Für das Laden der Parametersätze ist eine Schnittstelle zum Zugriff auf die LEA-interne Parameterdatenhaltung (Array) notwendig. Eine solche Schnittstelle ist im MTP-Konzept derzeit nicht vorgesehen und wird aufgrund dessen als \emph{ArrayServParam}-Schnittstelle neu spezifiziert. Die einzelnen Parametersätze des Arrays weisen, wie im Zuge der \emph{Übertragung einzelner Parametersätze} beschrieben, einen LEA-spezifischen strukturierten Datentyp auf. 

\subsection{Initiierung der Parametrierung}
\label{subsec:ParameterInitiierung}

Im Kontext der Parametrierung von LEA-Diensten lassen sich die in Tabelle~\ref{tab:ParametrierungsInitiierung} dargestellten Arten der Initiierung der Diensteparametrierung unterscheiden \cite{Blumenstein.atpLEA,Janzen.2023,Blumenstein.MoProLog}. Nachfolgend wird beschrieben, welche Vor- und Nachteile diese Mechanismen aufweisen und wie sie mit bestehenden und neuen MTP-Konzepten umgesetzt werden können.

% Bewertung Initiierung der Parametrierung
\begin{footnotesize}
	\begin{longtable}{|
		>{\columncolor[HTML]{E0E0E0}}l 
		>{\columncolor[HTML]{FFFFFF}}l |}
		\caption{Arten der Initiierung der Parametrierung von Logistics Equipment Assemblies}
		\label{tab:ParametrierungsInitiierung}\\
		\hline

		\multicolumn{2}{|L{14.8cm}|}{\cellcolor[HTML]{E0E0E0}\textbf{Parametrierung durch den LOL}} 
		\\ \hline

		\multicolumn{2}{|L{14.8cm}|}{\cellcolor[HTML]{FFFFFF}{Dieser Mechanismus sieht vor, dass der LOL das Wissen besitzt, wann welche Parameter an die LEAs zu übergeben sind, und die Parametrierung entsprechend initiiert.}} 
		\\ \hline

		\multicolumn{1}{|L{1.5cm}|}{\cellcolor[HTML]{FFFFFF} Vorteile }        	& 
		\multicolumn{1}{L{13.13cm}|}{\cellcolor[HTML]{FFFFFF} Der LOL besitzt das Wissen über alle Parameter und überträgt diese auch selbst an die LEAs. Somit wird die Komplexität der Parameterübertragung vom LOL auf die LEAs, v.~a. im Rahmen der auftragsorientierten Arbeitsweise, gering gehalten. } 
		\\ \hline

		\multicolumn{1}{|L{1.5cm}|}{\cellcolor[HTML]{FFFFFF} Nachteile }        	& 
		\multicolumn{1}{L{13.13cm}|}{\cellcolor[HTML]{FFFFFF} Der LOL muss stets den Status der (bedarfsorientiert arbeitenden) LEAs verfolgen und ermitteln, wann eine Parametrierung erforderlich ist. Im Falle der Parametrierung durch \emph{Übertragung einer Auswahl an Parametersätzen} muss der LOL zudem verfolgen, welche Parametersätze in den lokalen Datenhaltungen der LEAs bereits vorliegen.} 
		\\ \hline

		\pagebreak

		\hline
		\multicolumn{2}{|L{14.8cm}|}{\cellcolor[HTML]{E0E0E0}\textbf{Parameteranfrage durch die LEA}} 
		\\ \hline

		\multicolumn{2}{|L{14.8cm}|}{\cellcolor[HTML]{FFFFFF}{Dieser Mechanismus sieht vor, dass die LEA selbst feststellen kann, wann ihr welche Parameter fehlen und diese vom LOL anfordert.}} 
		\\ \hline

		\multicolumn{1}{|L{1.5cm}|}{\cellcolor[HTML]{FFFFFF} Vorteile }        	& 
		\multicolumn{1}{L{13.13cm}|}{\cellcolor[HTML]{FFFFFF} Die LEA kann selbstständig feststellen, wann sie einen Parametrierungsbedarf hat, und welche Parameter ihr fehlen. Somit wird die Komplexität der Meldung von Parametrierungsbedarfen von LEAs an den LOL, v.~a. im Rahmen der bedarfsorientierten Arbeitsweise, gering gehalten. } 
		\\ \hline

		\multicolumn{1}{|L{1.5cm}|}{\cellcolor[HTML]{FFFFFF} Nachteile }        	& 
		\multicolumn{1}{L{13.13cm}|}{\cellcolor[HTML]{FFFFFF} Die LEAs erkennen ihre Parametrierungsbedarfe selbstständig, die Verwaltung aller Parameterwerte erfolgt allerdings im Parameter-Management des LOL. Somit muss zunächst der Parametrierungsbedarf durch eine LEA an den LOL gemeldet werden, damit dieser anschließend die entsprechenden Parameter an die LEA überträgt. }
		\\ \hline

		\multicolumn{2}{|L{14.8cm}|}{\cellcolor[HTML]{E0E0E0}\textbf{Parametrierung am lokalen LEA-Bedienbild}} 
		\\ \hline

		\multicolumn{2}{|L{14.8cm}|}{\cellcolor[HTML]{FFFFFF}{Dieser Mechanismus sieht vor, dass Parameter durch einen Operator lokal am LEA-Bedienbild eingegeben werden. Dabei ist zu berücksichtigen, dass Parameteränderungen ggf. in den LOL zurückgelesen werden müssen.}} 
		\\ \hline

		\multicolumn{1}{|L{1.5cm}|}{\cellcolor[HTML]{FFFFFF} Vorteile }        	& 
		\multicolumn{1}{L{13.13cm}|}{\cellcolor[HTML]{FFFFFF} Der Operator besitzt das Wissen über die notwendigen Parameter und überträgt diese auch selbst an die LEAs. Somit wird die Komplexität der Parameterübertragung vom Operator auf die LEAs, v.~a. im Rahmen der auftragsorientierten Arbeitsweise, gering gehalten.} 
		\\ \hline

		\multicolumn{1}{|L{1.5cm}|}{\cellcolor[HTML]{FFFFFF} Nachteile }        	& 
		\multicolumn{1}{L{13.13cm}|}{\cellcolor[HTML]{FFFFFF} Der Operator muss stets den Status der (bedarfsorientiert arbeitenden) LEAs verfolgen und ermitteln, wann eine Parametrierung erforderlich ist. Im Falle der Parametrierung durch \emph{Übertragung einer Auswahl an Parametersätzen} muss der Operator zudem verfolgen, welche Parametersätze in den lokalen Datenhaltungen der LEAs bereits vorliegen.} 
		\\ \hline
		
	\end{longtable}
\end{footnotesize}

Diese Arten der Parametrierung sind gemäß REQ-LEA-1e und REQ-MLS-1 im Kontext der LEA-Automatisierung MTP-basiert umzusetzen.

\textbf{Parametrierung durch den LOL:} Dieser Mechanismus entspricht der Parametrierung, die in den bisherigen MTP-Konzepten vorgesehen ist. Der LOL kann die Parametrierung unter Nutzung der in \cite{PNO.2025.Part4} spezifizierten Parameterschnittstellen vornehmen. Es ist keine Erweiterung der MTP-Spezifikation notwendig.

\textbf{Parameteranfrage durch die LEA:} Ein solcher Mechanismus ist bisher in der MTP-Spezifikation nicht vorgesehen. Anfragen einer PEA an einen POL können gemäß \cite{PNO.2025.Part4} mittels des Konzepts der Dienst-Bediener-In\-ter\-ak\-ti\-on umgesetzt werden. Es werden daher Mechanismen zur Parameteranfrage als Erweiterung der Dienst-Bediener-Interaktion spezifiziert. Für den in Abschnitt~\ref{subsec:Parametrierngsmech} empfohlenen Parametrierungsmechanismus wird ein \emph{ProductParameterRequest} und ein \emph{PackagingParameterRequest} spezifiziert. Diese ermöglichen es einer LEA, passende produkt- und verpackungsspezifische Parameter beim LOL anzufragen. 

\textbf{Parametrierung am lokalen LEA-Bedienbild:} Eine Parametrierung am lokalen LEA-Bedienbild ist bereits in der MTP-Spezifikation vorgesehen. Diese erfolgt je nach Datentyp der Parameter mittels der in \cite{PNO.2025.Part4} spezifizierten Parameterschnittstellen oder der neu spezifizierten Schnittstellen \emph{StructServParam} und \emph{ArrayServParam} (siehe auch Abschnitt~\ref{subsec:DynHmiObjekte}). Eine Möglichkeit für LEAs, Parameteränderungen an den LOL zu melden, ist allerdings in der MTP-Spezifikation bisher nicht vorgesehen. Es werden daher Mechanismen zur Meldung von Parameteränderungen als Erweiterung der Dienst-Bediener-Interaktion spezifiziert. Für den in Abschnitt~\ref{subsec:Parametrierngsmech} empfohlenen Parametrierungsmechanismus wird ein \emph{ProductParameterUpdatedInfo} und ein \emph{PackagingParameterUpdatedInfo} spezifiziert. Diese ermöglichen es einer LEA, Parameteränderungen ihrer produkt- und verpackungsspezifischen Parameter beim LOL zu melden. 

\subsection{Empfohlener Parametrierungsmechanismus}
\label{subsec:Parametrierngsmech}

Die vorgestellten Übertragungsarten (Abschnitt~\ref{subsec:ParameterUebertragung}) und Initiierungsarten (Abschnitt~\ref{subsec:ParameterInitiierung}) stellen Grundbausteine für die LEA-Parametrierung dar, die grund\-sätz\-lich beliebig kombiniert werden können. In der studentischen Arbeit \cite{Janzen.2023} wurde darauf basierend eine Empfehlung erarbeitet, wie die in Abschnitt~\ref{subsec:Parameterarten} beschriebenen Parameterarten an LEA-Dienste übertragen werden sollten.

\subsubsection{Auftragsspezifische Parameter}

Auftragsspezifische Parameter werden in der Regel durch ein Auftrags-Management des LOL verwaltet und übertragen. Bei MLSs mit einer geringen Anzahl an LEAs können sie auch durch einen Operator an einem LEA-Bedienbild im LOL eingestellt werden. Da eine LEA in der Regel nur eine geringe Anzahl auftragsspezifischer Parameter erhält, werden diese häufig in Form von Einzelparametern über Parameterschnittstellen gemäß \cite{PNO.2025.Part4} an die LEA übertragen. Ist für eine LEA ein umfassenderer Satz auftragsspezifischer Parameter notwendig, so ist auch eine Übertragung als einzelner Parametersatz über eine \emph{StructServParam}-Schnittstelle möglich, um den Parametrierungsaufwand zu reduzieren. Aufgrund ihres auftragsbezogenen Charakters sind diese Parameter als Prozedurparameter gemäß \cite{PNO.2025.Part4} umzusetzen.

\subsubsection{Produkt- und verpackungsspezifische Parameter}

Die häufig sehr umfangreichen produkt- und verpackungsspezifischen Parameter werden durch ein Parameter-Management des LOL verwaltet. Dieses beinhaltet eine strukturierte Datenbank zur Verwaltung notwendiger Parameter für alle zu verpackenden Produkte und Verpackungsmaterialien. Unter Nutzung dieser Datenbasis werden LEA-Dienste nach Bedarf parametriert. Aufgrund der großen Anzahl wird für produkt- und verpackungsspezifische Parameter eine Parametrierung durch eine Auswahl von Parametersätzen empfohlen, um einen möglichst geringen Parametrierungs- und Kommunikationsaufwand zu erreichen. Abbildung~\ref{fig:ProdPackParametrierung} stellt den empfohlenen Parametrierungsmechanismus für produkt- und verpackungsspezifische LEA-Parameter dar, der nachfolgend beschrieben wird.

\begin{figure}[htbp]
	\centerline{\includegraphics[width=0.97\textwidth]{Inhalt/Abbildungen/06_LEA-Automatisierung/Parameter/Empfohlene_Parametrierung_LOL_LEA.pdf}}
	\caption{Empfohlener Parametrierungsmechanismus für produkt- und verpackungsspezifische LEA-Parameter}
	\label{fig:ProdPackParametrierung}
\end{figure} 

\textbf{Struktur der LEA-internen Parametersatzverwaltung:} Die Pa\-ra\-me\-ter\-sät\-ze werden in einer LEA-internen Datenverwaltung, z.~B. einem LEA-internen Datenbaustein, abgelegt. Die prinzipielle Struktur dieser Datenverwaltung ist in Abbildung~\ref{fig:LeaParameterverwaltung} dargestellt. 

\begin{figure}[htbp]
	\centerline{\includegraphics[width=0.9\textwidth]{Inhalt/Abbildungen/06_LEA-Automatisierung/Parameter/Product_PackagingDataSet.pdf}}
	\caption{Struktur der internen Parameterverwaltung einer LEA}
	\label{fig:LeaParameterverwaltung}
\end{figure}

Es werden zwei Arrays vorgesehen -- ein \emph{ProductDataSet} für die produktspezifischen und ein \emph{PackagingDataSet} für die verpackungsspezifischen Parameter.\footnote{Grundsätzlich könnten auch alle produkt- und verpackungsspezifischen Parameter in einem einzigen Array (\emph{ProductDataSet}) abgelegt werden. Dies führt zwar zu einem einfacheren Parametrierungsmechanismus, aber auch zu einer ggf. redundanten Speicherung der verpackungsspezifischen Parameter in den produktspezifischen Parametersätzen. Zudem können verpackungsspezifische Parameter nicht unabhängig von den produktspezifischen Parametern abgerufen werden. Aufgrund dessen wurde die Parameterverwaltung in zwei getrennten Arrays gewählt.} Die einzelnen Array-Elemente weisen jeweils einen LEA-spezifischen strukturierten Datentyp auf. Ein Parametersatz im \emph{ProductDataSet} kann mittels einer \emph{ProductId} und eines \emph{LogisticsObjectStatus} als zwei Schlüssel eindeutig referenziert werden. Die \emph{ProductId} ist eine eindeutige Kennung für das zu verpackende Produkt, der \emph{LogisticsObjectStatus} beschreibt den aktuellen Bearbeitungszustand des LO. Ein Parametersatz im \emph{PackagingDataSet} kann mittels einer \emph{PackagingId}, welche eine eindeutige Kennung für Verpackungsmaterialien darstellt, referenziert werden. Wird für ein Produkt ein bestimmtes Verpackungsmaterial benötigt, wird dessen \emph{PackagingId} im Parametersatz des Produkts hinterlegt, was eine Referenz vom \emph{ProductDataSet} auf das \emph{PackagingDataSet} darstellt.

\textbf{Parametrierung auf Initiative des LOL oder der LEA:} Das \emph{ProductDataSet} und das \emph{PackagingDataSet} werden mittels je einer \emph{ArrayServParam}-Schnittstelle an der LEA befüllt, über die Parametersätze an die LEA übergeben und an definierten Stellen des jeweiligen Arrays hinterlegt werden (Abbildung~\ref{fig:LeaParameterverwaltung}). Diese Befüllung der Arrays kann zu jeder Zeit (unabhängig von anstehenden Verpackungsaufträgen) durch das Parameter-Management des LOL durchgeführt werden. \\
Ist für die anstehende Bearbeitung eines LO \underline{kein} passender produkt- oder verpackungsspezifischer Parametersatz in der LEA verfügbar, kann die LEA diesen über einen \emph{ProductParameterRequest} bzw. \emph{PackagingParameterRequest} anfragen (Abbildung~\ref{fig:ProdPackParametrierung}). Einem \emph{ProductParameterRequest} gibt die LEA die \emph{ProductId} und den \emph{LogisticsObjectStatus} des gewünschten produktspezifischen Parametersatzes mit. Einem \emph{PackagingParameterRequest} gibt die LEA die \emph{PackagingId} des gewünschten verpackungsspezifischen Parametersatzes mit. Das Parameter-Management schreibt anschließend den gewünschten Parametersatz an die entsprechende \emph{ArrayServParam}-Schnittstelle der LEA. \\
Somit kann eine Vorbefüllung der Arrays durch den LOL oder eine Parameteranfrage durch die LEA erfolgen. Eine Empfehlung, wann welche der beiden Varianten zu bevorzugen ist, ist in Tabelle~\ref{tab:VorbefüllungArrays} zu finden. 

% Bewertung Vorbefüllung Arrays
\begin{footnotesize}
	\begin{longtable}{|
		>{\columncolor[HTML]{E0E0E0}}l 
		>{\columncolor[HTML]{FFFFFF}}l |}
		\caption{Empfehlung zur Vorbefüllung des \emph{ProductDataSet} und des \emph{PackagingDataSet}}
		\label{tab:VorbefüllungArrays}\\
		\hline

		\multicolumn{2}{|L{14.8cm}|}{\cellcolor[HTML]{E0E0E0}\textbf{Vorbefüllen der Arrays}} 
		\\ \hline

		\multicolumn{2}{|L{14.8cm}|}{\cellcolor[HTML]{FFFFFF}Dieser Mechanismus sieht vor, dass das \emph{ProductDataSet} und das \emph{PackagingDataSet} vor der Übergabe von Aufträgen an die LEA (z.~B. während der Inbetriebnahme der LEA) mit allen für die Bearbeitung notwendigen Parametersätzen befüllt werden. } 
		\\ \hline

		\multicolumn{1}{|L{2.3cm}|}{\cellcolor[HTML]{FFFFFF} Anwendungs-fall }        	& 
		\multicolumn{1}{L{12.33cm}|}{\cellcolor[HTML]{FFFFFF} Diese Variante wird bei einem festen zu verpackenden Produktportfolio empfohlen, welches sich selten ändert.} 
		\\ \hline

		\multicolumn{1}{|L{2.3cm}|}{\cellcolor[HTML]{FFFFFF} Vorteile }        	& 
		\multicolumn{1}{L{12.33cm}|}{\cellcolor[HTML]{FFFFFF} Solange keine neuen Parametersätze in die Arrays übertragen werden müssen, ist keine Kommunikation zwischen den LEAs und dem Parameter-Management im LOL zur Laufzeit notwendig. Die Netzwerklast ist gering und die LEA-Dienste können u.~U. auch ohne Kommunikationsverbindung zum Parameter-Management (z.~B. auch bei Ausfall des Parameter-Managements) ausgeführt werden.} 
		\\ \hline

		\multicolumn{2}{|L{14.8cm}|}{\cellcolor[HTML]{E0E0E0}\textbf{Anfragen der Parametersätze zur Laufzeit}} 
		\\ \hline

		\multicolumn{2}{|L{14.8cm}|}{\cellcolor[HTML]{FFFFFF}Dieser Mechanismus sieht vor, dass das \emph{ProductDataSet} und das \emph{PackagingDataSet} bei Inbetriebnahme der LEA leer gelassen werden und die LEA notwendige Parametersätze bei Auftragserhalt anfragt. } 
		\\ \hline

		\multicolumn{1}{|L{2.3cm}|}{\cellcolor[HTML]{FFFFFF} Anwendungs-fall }        	& 
		\multicolumn{1}{L{12.33cm}|}{\cellcolor[HTML]{FFFFFF} Diese Variante wird bei einem sich häufig ändernden oder sehr großen Produktportfolio empfohlen, oder wenn eine LEA schnell in Betrieb genommen werden muss. }
		\\ \hline 

		\multicolumn{1}{|L{2.3cm}|}{\cellcolor[HTML]{FFFFFF} Vorteile }        	& 
		\multicolumn{1}{L{12.33cm}|}{\cellcolor[HTML]{FFFFFF} Die Arrays werden nach Bedarf zur Laufzeit befüllt. Nicht mehr benötigte Parametersätze können gelöscht oder überschrieben werden. Eine LEA kann schnell in Betrieb genommen werden und sofort einen Verpackungsauftrag entgegennehmen, ohne dass erst alle Parametersätze vorab geladen werden müssen.} 
		\\ \hline

		\multicolumn{2}{|L{14.8cm}|}{\cellcolor[HTML]{E0E0E0}\textbf{Weiterführende Empfehlung}} 
		\\ \hline

		\multicolumn{2}{|L{14.8cm}|}{\cellcolor[HTML]{FFFFFF} Unabhängig davon, mit welcher Variante die Arrays befüllt werden, sollte sichergestellt werden, dass für jeden im \emph{ProductDataSet} hinterlegten produktspezifischen Parametersatz der darin referenzierte verpackungsspezifische Parametersatz im \emph{PackagingDataSet} hinterlegt ist. Wenn ein neuer produktspezifischer Parametersatz an eine LEA übertragen wird, wird daher empfohlen, immer zu prüfen, ob der darin referenzierte verpackungsspezifische Parametersatz bereits in der LEA hinterlegt ist. Ist dies nicht der Fall, sollte dieser auch an die LEA übertragen werden. Diese Prüfung kann im Parameter-Management des LOL oder LEA-intern durchgeführt werden. Entsprechend erfolgt die ggf. notwendige Übertragung des verpackungsspezifischen Parametersatzes entweder auf Initiative des LOL oder auf Anfrage der LEA.} 
		\\ \hline

	\end{longtable}
\end{footnotesize}

\pagebreak

Da das Laden des \emph{ProductDataSet} und des \emph{PackagingDataSet} unabhängig von einem konkreten Verpackungsauftrag durchgeführt werden kann, sind die zugehörigen \emph{ArrayServParam}-Schnittstellen als Konfigurationsparameter gemäß \cite{PNO.2025.Part4} auszuführen.
 
\textbf{Auswahl geladener Parametersätze:} Für die Bearbeitung eines LO sind die passenden Parametersätze aus beiden Arrays auszuwählen. Dies erfolgt unterschiedlich für Dienste im CES- und im SES-Betrieb. \textbf{CES-Betrieb:} Bei einem Dienst im CES-Betrieb wird die \emph{ProductId} über einen Prozedurparameter vom Typ DINT (Schnittstellendefinition \emph{DIntServParam} gemäß \cite{PNO.2025.Part4}) beim Starten eines Auftrags als auftragsspezifischer Einzelparameter vorgegeben. Da im CES-Betrieb nicht auf unterschiedliche Bearbeitungsstatus von LOs reagiert werden muss, ist die Vorgabe des \emph{LogisticsObjectStatus} nicht zwingend notwendig. Bei Bedarf kann dieser durch einen Prozedurparameter vom Typ DINT eingestellt werden, um den Bearbeitungsstatus des LO explizit vorzugeben. Entsprechend der Vorgabe wird der passende produktspezifische Parametersatz ausgewählt. \textbf{SES-Betrieb:} Bei einem Dienst im SES-Betrieb werden die \emph{ProductId} und der \emph{LogisticsObjectStatus} als Teil der Transportauftragsdaten (Abschnitt~\ref{subsec:SchnittstelleTransportdienst}) beim Eintreffen eines LO an die LEA übermittelt, somit sind keine separaten Parameterschnittstellen erforderlich. Ist der produktspezifische Parametersatz ausgewählt, wird im CES- und auch im SES-Betrieb die \emph{PackagingId} ermittelt und der zugehörige verpackungsspezifische Parametersatz ausgewählt.

\textbf{Parametrierung am lokalen LEA-Bedienbild:} Während der LEA-In\-be\-trieb\-nah\-me ergeben sich häufig noch Änderungen an produkt- oder verpackungsspezifischen Parametersätzen, die lokal am LEA-Bedienbild vorgenommen und vor Ort getestet werden. Das Parameter-Management des LOL wird dann mittels \emph{ProductParameterUpdatedInfo} bzw. \emph{PackagingParameterUpdatedInfo} unter Angabe eines Arrayindex informiert, dass sich der Parametersatz an entsprechender Stelle des \emph{ProductDataSet} bzw. \emph{PackagingDataSet} geändert hat. Der LOL entscheidet, ob diese Änderung in die Datenhaltung des Parameter-Managements übernommen werden soll, um sie z.~B. auch auf andere LEAs auszurollen. In jedem Fall bestätigt das Parameter-Management dem LEA-Dienst, dass die Änderung zur Kenntnis genommen wurde. Abbildung~\ref{fig:HmiParametrierung} zeigt diese Meldung zur Parametersatzänderung einer LEA an den LOL. 

\begin{figure}[htb]
	\centerline{\includegraphics[width=0.85\textwidth]{Inhalt/Abbildungen/06_LEA-Automatisierung/Parameter/Empfohlene_Parametrierung_HMI.pdf}}
	\caption{Meldung einer LEA an den LOL bei Änderungen an produkt- und verpackungsspezifischen LEA-Parametersätzen}
	\label{fig:HmiParametrierung}
\end{figure} 

\pagebreak
 
\textbf{Semantik der Parameterschnittstellen:} Zur Umsetzung der beschriebenen Mechanismen muss das Parameter-Management des LOL eindeutig identifizieren können, welche \emph{ArrayServParam}-Schnittstellen zur Verwaltung des \emph{ProductDataSet} und des \emph{PackagingDataSet} verwendet werden. Zudem müssen die \emph{DIntServParam}-Schnittstelle für die Eingaben der \emph{ProductId} und des \emph{LogisticsObjectStatus} für Dienste im CES-Betrieb eindeutig identifizierbar sein. Zu diesem Zweck werden \emph{FunctionClassificationAttributes}, die in der bisherigen MTP-Spezifikation nur für Dienste und Prozeduren vorgesehen sind, auch für Parameter eingeführt. Weiterhin werden konkrete \emph{FunctionClassificationAttributes} für die Parameter \emph{ProductId}, \emph{ProductDataSet} und \emph{PackagingDataSet} spezifiziert. 

\subsubsection{Konstruktionsspezifische Parameter}

Konstruktionsspezifische Parameter werden meist zum Inbetriebnahmezeitpunkt durch einen Operator an die LEA übergeben. Diese Parameter betreffen den physischen Aufbau einer LEA (z.~B. Vorgabe des verbauten Abfüllstutzens in einer Abfüll-LEA) bzw. die Bestückung einer LEA (z.~B. eingelegter Palettentyp in einer Palettieren-LEA). Sie sind unabhängig von einem konkreten Auftrag und unabhängig davon, ob eine LEA im CES- oder SES-Betrieb betrieben wird. Es werden stattdessen grundlegende Einstellungen für die Diensteausführung vorgenommen. Somit sind diese Parameter als Konfigurationsparameter gemäß \cite{PNO.2025.Part4} umzusetzen. Die Anzahl der konstruktionsspezifischen Parameter ist meist gering. Daher werden sie in der Regel als Einzelwerte über bestehende Schnittstellen gemäß \cite{PNO.2025.Part4} an die LEA übertragen.

Bei LEA-Typen, die ein bestimmtes Verpackungsmaterial nutzen, ist mittels eines konstruktionsspezifischen Einzelparameters vom Typ \emph{DIntServParam} gemäß \cite{PNO.2025.Part4} einzustellen, welches Verpackungsmaterial (z.~B. Palettentyp) aktuell in der LEA \emph{vorhanden ist}. Auf Grundlage dieses Parameters können entsprechende verpackungsspezifische Informationen aus dem LEA-internen \emph{PackagingDataSet} abgerufen oder beim LOL angefragt werden (ähnlich Abbildung~\ref{fig:ProdPackParametrierung}, unterer Teil). Dieser Wert gleicht somit semantisch der im vorherigen Abschnitt beschriebenen \emph{PackagingId} und ist vor der Bearbeitung eines LO mit der \emph{PackagingId} abzugleichen, die im gewählten produktspezifischen Parametersatz \emph{gefordert ist}. Passen diese nicht überein, kann die LEA den Auftrag nicht bearbeiten oder muss mit einem anderen Verpackungsmaterial bestückt werden. Um den \emph{Packaging\-Id}-Konfigurationsparameter semantisch eindeutig identifizieren zu können, wird ein entsprechendes \emph{FunctionClassificationAttribute} eingeführt. 

\section{Reportwerte}
\label{sec:Reportwerte}

Das in \cite{PNO.2025.Part4} beschriebene Konzept für Reportwerte ist auch für den Einsatz im Bereich der produktionsnahen Logistik geeignet. Erweiterungen sind lediglich in Bezug auf die verfügbaren Datentypen notwendig. 

Es ist sinnvoll, für die im vorhergehenden Abschnitt~\ref{sec:LeaParameter} eingeführten strukturierten Datentypen und Arrays auch entsprechende Reportwerte zu spezifizieren, um beispielsweise die aktuell in der LEA vorhandenen Parametersätze dokumentieren zu können. Gemäß \cite{PNO.2025.Part4} werden als Grundlage für die Bereitstellung und Modellierung von Reportwerten die in \cite{PNO.2025.Part3} spezifizierten Ableitungen der \emph{SUC IndicatorElement} verwendet. Mit \emph{StructView} und \emph{ArrayView} werden entsprechend zwei neue Schnittstellendefinitionen eingeführt, die ebenso diesem Prinzip des MTP-Konzepts folgen. 

Gemäß \cite{PNO.2025.Part4} können Reportwerte \enquote{eingefroren} werden, d.~h. der aktuelle Wert kann für eine gewisse Zeit beibehalten werden. Optional kann zudem \emph{MissedValueFlag} ergänzt werden, das anzeigt, ob sich der Reportwert mindestens zweimal geändert hat, während er eingefroren war. Diese Prinzipien werden auch für Reportwerte mit strukturierten Datentypen und Array-Datentypen übernommen. Bei Array-Datentypen sollten im Falle des Einfrierens alle Werte des Arrays eingefroren werden und angezeigt werden können.

\section{Prozesswerte}
\label{sec:Prozesswerte}

Das in \cite{PNO.2025.Part4} beschriebene Konzept für Prozesswerte ist auch für den Einsatz im Bereich der produktionsnahen Logistik geeignet. Erweiterungen am grundlegenden Konzept sind nicht notwendig. Im Gegensatz zu Parametern und Reportwerten wurde im Zuge dieser Arbeit \underline{kein} Bedarf zur Übertragung von Prozesswerten mit strukturierten Datentypen oder Array-Datentypen bei der Automatisierung von MLSs identifiziert. Da im MTP-Konzept allerdings zu jedem Datentyp auch immer Prozesswerte vorgesehen sind, werden entsprechende Schnittstellendefinitionen zur Übertragung von Prozesswerten mit strukturierten Datentypen und Array-Datentypen spezifiziert. Im Speziellen umfasst das die Schnittstellendefinitionen \emph{StructProcessValueIn}, \emph{ArrayProcessValueIn} und \emph{ArrayProcessValueOut}. Der Prozesswertausgang für strukturierte Datentypen entspricht der Schnittstellendefinition \emph{StructView} und muss nicht separat spezifiziert werden.

\section{Bedienbilder}
\label{sec:LeaHmi}

Zur Visualisierung der ablaufenden Logistikprozesse steht im LOL eine Bedienbild-Funktionalität bereit, über die die LEAs überwacht und manuell gesteuert werden können. In diese sind die Bedienbilder einzelner LEAs aufbauend auf den Konzepten aus \cite{PNO.2025.Part2} automatisiert zu integrieren. Während die Konzepte aus \cite{PNO.2025.Part2} für R\&I\footnote{\gls{ri} sind detaillierte Diagramme zur Abbildung verfahrenstechnischer Anlagen. Darin werden mittels einer einheitlichen Symbolik, Anlagenteile, Rohrleitungen sowie Mess-, Steuer- und Regelfunktionen abgebildet. \cite{DINNormenausschussChemischerApparatebau.2015}}-ähnliche Bedienbilder konzipiert sind, sind für den Bereich der produktionsnahen Logistik maschinenorientierte Bedienbilder, wie in Abbildung~\ref{fig:PalHmi} für einen Palettierer dargestellt, üblich. Diese Bedienbilder enthalten ein Abbild der LEA als statisches Bedienbildobjekt und mehrere dynamische Bedienbildobjekte für den LEA-Dienst, Parameter, Reportwerte, Prozesswerte und weitere LEA-interne Werte. 

\begin{figure}[htbp]
	\centerline{\includegraphics[width=1.0\textwidth]{Inhalt/Abbildungen/06_LEA-Automatisierung/HMI/Example_HMI_PAL.png}}
	\caption{Bedienbild eines Palettierers}
	\label{fig:PalHmi}
\end{figure}

\subsection{Dynamische Bedienbildobjekte}
\label{subsec:DynHmiObjekte}

Dynamische Bedienbildobjekte können grundlegend mit den Mechanismen aus \cite{PNO.2025.Part2} implementiert werden. Für die Anzeige und Manipulation verschiedener LEA-interner Variablen werden die Schnittstellendefinitionen gemäß \cite{PNO.2025.Part3} und \cite{PNO.2025.Part4} genutzt. Zudem werden im Kontext der LEA-Automatisierung Schnittstellendefinitionen für die Anzeige und Manipulation von Variablen mit strukturierten Datentypen und Array-Datentypen ergänzt. Tabelle~\ref{tab:DynamischeHmiObjekte} fasst diese neu eingeführten Schnittstellen zusammen, die als Grundlage für die dynamischen Bedienbildobjekte dienen. Sie orientieren sich an den Schnittstellen, die die MTP-Spezifikation bereits für primitive Datentypen vorsieht.

% Verwendung der Schnittstellendefinitionen für die Bedienbildmodellierung
\begin{footnotesize}
	\begin{longtable}{|
		>{\columncolor[HTML]{E0E0E0}}l 
		>{\columncolor[HTML]{FFFFFF}}l |}
		\caption{Schnittstellendefinitionen für die Modellierung dynamischer Bedienbildobjekte}
		\label{tab:DynamischeHmiObjekte}\\
		\hline

		\multicolumn{1}{|L{3.5cm}|}{\cellcolor[HTML]{E0E0E0} \textbf{Schnittstelle}}        	& 
		\multicolumn{1}{L{11.13cm}|}{\cellcolor[HTML]{E0E0E0}\textbf{Verwendung}} 
		\\ \hline

		\multicolumn{1}{|L{3.5cm}|}{\cellcolor[HTML]{FFFFFF}  StructView}        	& 
		\multicolumn{1}{L{11.13cm}|}{\cellcolor[HTML]{FFFFFF} Anzeige von (Report-) Werten und Prozesswertausgängen mit strukturierten Datentypen} 
		\\ \hline

		\multicolumn{1}{|L{3.5cm}|}{\cellcolor[HTML]{FFFFFF}  ArrayView}        	& 
		\multicolumn{1}{L{11.13cm}|}{\cellcolor[HTML]{FFFFFF} Anzeige von (Report-) Werten mit Array-Datentypen} 
		\\ \hline

		\multicolumn{1}{|L{3.5cm}|}{\cellcolor[HTML]{FFFFFF} StructMan }        	& 
		\multicolumn{1}{L{11.13cm}|}{\cellcolor[HTML]{FFFFFF} Manipulation von Werten strukturierter Datentypen durch einen Operator} 
		\\ \hline

		\multicolumn{1}{|L{3.5cm}|}{\cellcolor[HTML]{FFFFFF} StructManInt }        	& 
		\multicolumn{1}{L{11.13cm}|}{\cellcolor[HTML]{FFFFFF} Manipulation von Werten strukturierter Datentypen durch einen Operator oder die LEA selbst} 
		\\ \hline

		\multicolumn{1}{|L{3.5cm}|}{\cellcolor[HTML]{FFFFFF} ArrayMan }        	& 
		\multicolumn{1}{L{11.13cm}|}{\cellcolor[HTML]{FFFFFF} Manipulation von Werten mit Array-Datentypen durch einen Operator} 
		\\ \hline
		
		\multicolumn{1}{|L{3.5cm}|}{\cellcolor[HTML]{FFFFFF} ArrayManInt }        	& 
		\multicolumn{1}{L{11.13cm}|}{\cellcolor[HTML]{FFFFFF} Manipulation von Werten mit Array-Datentypen durch einen Operator oder die LEA selbst} 
		\\ \hline

		\multicolumn{1}{|L{3.5cm}|}{\cellcolor[HTML]{FFFFFF} StructServParam }        	& 
		\multicolumn{1}{L{11.13cm}|}{\cellcolor[HTML]{FFFFFF} Anzeige und Manipulation von Parametern mit strukturierten Datentypen} 
		\\ \hline

		\multicolumn{1}{|L{3.5cm}|}{\cellcolor[HTML]{FFFFFF} ArrayServParam }        	& 
		\multicolumn{1}{L{11.13cm}|}{\cellcolor[HTML]{FFFFFF} Anzeige und Manipulation von Parametern mit Array-Datentypen} 
		\\ \hline

		\multicolumn{1}{|L{3.5cm}|}{\cellcolor[HTML]{FFFFFF}  StructProcessValueIn }        	& 
		\multicolumn{1}{L{11.13cm}|}{\cellcolor[HTML]{FFFFFF} Anzeige von Prozesswerteingängen mit strukturierten Datentypen } 
		\\ \hline

		\multicolumn{1}{|L{3.5cm}|}{\cellcolor[HTML]{FFFFFF}  ArrayProcessValueIn }        	& 
		\multicolumn{1}{L{11.13cm}|}{\cellcolor[HTML]{FFFFFF} Anzeige von Prozesswerteingängen mit Array-Datentypen } 
		\\ \hline

		\multicolumn{1}{|L{3.5cm}|}{\cellcolor[HTML]{FFFFFF}  ArrayProcessValueOut }        	& 
		\multicolumn{1}{L{11.13cm}|}{\cellcolor[HTML]{FFFFFF} Anzeige von Prozesswertausgängen mit Array-Datentypen } 
		\\ \hline
		
		
	\end{longtable}
\end{footnotesize}

Die Schnittstellen \emph{StructView} und \emph{ArrayView} dienen der Anzeige von LEA-in\-ter\-nen Werten. Wie in den Abschnitten~\ref{sec:Reportwerte} und \ref{sec:Prozesswerte} bereits eingeführt, können sie zur Anzeige von Reportwerten und im Falle von \emph{StructView} auch zur Anzeige von Prozesswertausgängen verwendet werden.

Zudem sieht das MTP-Konzept für jeden Datentyp \emph{*Mon}-Schnittstellen vor, die über die \emph{*View}-Schnittstellen hinaus Überwachungsfunktionen bereitstellen, ob ein Wert ober- oder unterhalb definierter Grenzen liegt. \emph{*Mon}-Schnittstellen werden für strukturierte Datentypen und Array-Datentypen \underline{nicht} spezifiziert, da diese Datentypen stets mehr als einen Einzelwert enthalten und somit die zuvor genannte Überwachung nicht sinnvoll ist.

\pagebreak

Darüber hinaus sieht das MTP-Konzept \emph{*Man}- und \emph{*ManInt}-Schnittstellen vor, die zur Manipulation von LEA-internen Werten durch einen Operator dienen. Die \emph{*ManInt}-Schnittstellen ermöglichen im Gegensatz zu den \emph{*Man}-Schnittstellen zu\-sätz\-lich die Manipulation der Werte durch die LEA selbst. Für strukturierte und arraybasierte Datentypen werden entsprechend die Schnittstellendefinitionen \emph{StructMan}, \emph{StructManInt}, \emph{ArrayMan} und \emph{ArrayManInt} spezifiziert.

Die Schnittstellen \emph{StructServParam} und \emph{ArrayServParam} dienen wie in Abschnitt~\ref{subsec:ParameterUebertragung} beschrieben der Anzeige und Manipulation von LEA-Diensteparametern mit strukturierten und arraybasierten Datentypen. 

Schließlich dienen die Schnittstellen \emph{StructProcessValueIn}, \emph{ArrayProcessValueIn} und \emph{ArrayProcessValueOut} wie in Abschnitt~\ref{sec:Prozesswerte} beschrieben zur Anzeige von Prozesswerteingängen mit strukturierten und arraybasierten Datentypen bzw. zur Anzeige von Prozesswertausgängen mit arraybasierten Datentypen.

Die beiden nachfolgenden Artefakte zur Automatisierung von Logistics Lines (Kapitel~\ref{chap:Art2LL}) und zur Automatisierung flexibler Transporte in Logistics Areas (Kapitel~\ref{chap:Art3LA}) stellen weitere Schnittstellen vor, für die ebenfalls dynamische Bedienbildobjekte in LEA-Bedienbildern vorgesehen werden können.

\subsection{Statische Bedienbildobjekte}
\label{subsec:StatHmiObjekte}

Statische Bedienbildobjekte werden gemäß \cite{PNO.2025.Part2} mittels \emph{VisualObjects} modelliert, die im Bedienbild positioniert und mit einer ECLASS-Referenz\footnote{ECLASS ist ein branchenübergreifendes Klassifizierungssystem, welches semantische Referenzen für Produkte und Dienstleistungen enthält. In \cite{ECLASSe.V..2025} kann die ECLASS-Datenbank nach passenden Referenzen durchsucht werden.} versehen werden. Das integrierende System (hier: LOL) muss diese Referenz in seiner Grafikbibliothek zur Verfügung haben, um das statische Objekt darstellen zu können. Im Falle von LEAs sollten die verwendeten Grafiken das LEA-Aussehen exakt nachbilden, um eine visuelle Beziehung zwischen dem Bedienbild und der realen LEA herstellen zu können. Da ein LOL nicht Bilder für alle denkbaren LEAs vorhalten kann, wird der sogenannte \emph{CustomSymbols}-Mechanismus zur automatisierten Integration von statischen LEA-Bildern (\emph{CustomSymbols}) eingeführt, der in Abbildung~\ref{fig:CustomSymbolHmi} dargestellt ist.

\begin{figure}[htbp]
	\centerline{\includegraphics[width=1.0\textwidth]{Inhalt/Abbildungen/06_LEA-Automatisierung/HMI/HMI_StatischesBild.drawio.pdf}}
	\caption{LEA-Bild als CustomSymbol im Anhang eines Palettierer-MTP}
	\label{fig:CustomSymbolHmi}
\end{figure}

CustomSymbols werden in Form von .svg-Dateien dem LEA-MTP als Anhänge gemäß \cite{PNO.2025.Part1} angefügt. Sie sind in einer separaten \emph{AttachmentGroup} gemäß \cite{PNO.2025.Part1} mit der Bezeichnung \enquote{CustomSymbols} zu organisieren. Zur Ablage der Datei im MTP-Anhang sollte aus Gründen der Übersichtlichkeit ein Ordner \gls{hmi} vorgesehen werden. 

Aus Sicht der Modellierung werden CustomSymbols genauso behandelt wie Grafiken aus der LOL-Grafikbibliothek. Zur Referenzierung eines solchen Bildes im MTP-Anhang wird dessen Dateiname auf die zugehörige ECLASS-Referenz (8-stelliger numerischer String) gesetzt, die auch bei der Bedienbildmodellierung in den \emph{VisualObjects} und für die Referenzierung mittels \emph{AttachmentReference} verwendet wird. Hierbei werden ECLASS-Nummern im Nummernbereich 9090XXXX verwendet, da diese nicht mit einer offiziellen Codierung belegt sind. Wird eine solche ECLASS in der Bedienbildmodellierung referenziert, verwendet der LOL das entsprechende CustomSymbol aus dem MTP-Anhang.

Diese Anwendung bestehender MTP-Mechanismen als Ergebnis dieser Dissertation wurde bereits in das Profil \emph{ModuleTypePackage:HMISet.Base V2.0.0} von \cite{PNO.2025.Part2} übernommen.

\section{MTP-Erweiterungen}
\label{sec:MTPErweiterungenLEA}

In den vorangegangenen Abschnitten wurde eine Reihe von Erweiterungen der MTP-Spezifikation identifiziert, die im Kontext der LEA-Automatisierung notwendig sind. Tabelle~\ref{tab:MtpErweiterungenLea} gibt einen Überblick über die hierfür einzuführenden und zu erweiternden Modell- und Schnittstellendefinitionen. Eine detailliertere Beschreibung erfolgt in den in Tabelle~\ref{tab:MtpErweiterungenLea} angegebenen Anhängen.

% MTP-Erweiterungen LEA
\begin{footnotesize}
	\begin{longtable}{|
		>{\columncolor[HTML]{E0E0E0}}l 
		>{\columncolor[HTML]{FFFFFF}}l 
		>{\columncolor[HTML]{FFFFFF}}l 
		>{\columncolor[HTML]{FFFFFF}}l |}
		\caption{Erweiterungen der MTP-Spezifikation für die Automatisierung von Logistics Equipment Assemblies}
		\label{tab:MtpErweiterungenLea}\\
		\hline

		% Schnittstellendefinitionen
		\multicolumn{4}{|L{14.8cm}|}{\cellcolor[HTML]{E0E0E0}\textbf{Schnittstellendefinitionen}} 
		\\ \hline

		\multicolumn{1}{|L{3.6cm}|}{\cellcolor[HTML]{E0E0E0}\textbf{Profil}}        	& 
		\multicolumn{1}{L{2.8cm}|}{\cellcolor[HTML]{E0E0E0}\textbf{Definition}} 			& 
		\multicolumn{1}{L{5.43cm}|}{\cellcolor[HTML]{E0E0E0}\textbf{Beschreibung}} 			& 
		\multicolumn{1}{L{2.45cm}|}{\cellcolor[HTML]{E0E0E0}\textbf{Detail-beschreibung}} 
		\\ \hline

		\multicolumn{1}{|L{3.6cm}|}{\cellcolor[HTML]{FFFFFF} ModuleTypePackage: ServiceSet.Complex-Types V2.0.0 }                    & 
		\multicolumn{1}{L{2.8cm}|}{\cellcolor[HTML]{FFFFFF} SUC Struct\-Serv\-Param}              &      
		\multicolumn{1}{L{5.43cm}|}{\cellcolor[HTML]{FFFFFF} Schnittstelle zur Übertragung von Dienste-Parametern strukturierter Datentypen}              &     
		\multicolumn{1}{L{2.45cm}|}{\cellcolor[HTML]{FFFFFF} Anhang~\ref{subsec:AnhangServiceSet}}             
		\\ \hline

		\multicolumn{1}{|L{3.6cm}|}{\cellcolor[HTML]{FFFFFF} ModuleTypePackage: ServiceSet.Complex-Types V2.0.0 }                    & 
		\multicolumn{1}{L{2.8cm}|}{\cellcolor[HTML]{FFFFFF} SUC Array\-Serv\-Param}              &      
		\multicolumn{1}{L{5.43cm}|}{\cellcolor[HTML]{FFFFFF} Schnittstelle zur Übertragung von Dienste-Parametern mit Array-Datentyp}              &     
		\multicolumn{1}{L{2.45cm}|}{\cellcolor[HTML]{FFFFFF} Anhang~\ref{subsec:AnhangServiceSet}}             
		\\ \hline

		\multicolumn{1}{|L{3.6cm}|}{\cellcolor[HTML]{FFFFFF} ModuleTypePackage: ServiceSet.Logistics V2.0.0 }         & 
		\multicolumn{1}{L{2.8cm}|}{\cellcolor[HTML]{FFFFFF} RC Logistics\-Interaction\-Extension}  & 
		\multicolumn{1}{L{5.43cm}|}{\cellcolor[HTML]{FFFFFF} Erweiterung der ServiceControl-Schnittstelle um die Variablen für LogisticsInteractions}              &     
		\multicolumn{1}{L{2.45cm}|}{\cellcolor[HTML]{FFFFFF} Anhang~\ref{subsec:AnhangServiceSet}}             
		\\ \hline

		\multicolumn{1}{|L{3.6cm}|}{\cellcolor[HTML]{FFFFFF} ModuleTypePackage: ServiceSet.Logistics V2.0.0 }                    & 
		\multicolumn{1}{L{2.8cm}|}{\cellcolor[HTML]{FFFFFF} SUC Service\-Control\newline (Erweiterung)}              &      
		\multicolumn{1}{L{5.43cm}|}{\cellcolor[HTML]{FFFFFF} Erweiterung der bestehenden ServiceControl-Schnittstelle um die Möglichkeit, die RC Logistics\-Interaction\-Extension anzufügen}  &     
		\multicolumn{1}{L{2.45cm}|}{\cellcolor[HTML]{FFFFFF} Anhang~\ref{subsec:AnhangServiceSet}}             
		\\ \hline

		\multicolumn{1}{|L{3.6cm}|}{\cellcolor[HTML]{FFFFFF} ModuleTypePackage: DataAssemblySet. ComplexTypes V2.0.0 }         & 
		\multicolumn{1}{L{2.8cm}|}{\cellcolor[HTML]{FFFFFF} SUC Struct\-View}              &      
		\multicolumn{1}{L{5.43cm}|}{\cellcolor[HTML]{FFFFFF} Schnittstelle zur Anzeige strukturierter Werte einer LEA}              &     
		\multicolumn{1}{L{2.45cm}|}{\cellcolor[HTML]{FFFFFF} Anhang~\ref{subsec:AnhangDataAssemblySet}}             
		\\ \hline
 
		\multicolumn{1}{|L{3.6cm}|}{\cellcolor[HTML]{FFFFFF} ModuleTypePackage: DataAssemblySet. ComplexTypes V2.0.0 }         & 
		\multicolumn{1}{L{2.8cm}|}{\cellcolor[HTML]{FFFFFF} SUC Array\-View}              &      
		\multicolumn{1}{L{5.43cm}|}{\cellcolor[HTML]{FFFFFF} Schnittstelle zur Anzeige von Werten eines LEA-intern verwalteten Arrays}              &     
		\multicolumn{1}{L{2.45cm}|}{\cellcolor[HTML]{FFFFFF} Anhang~\ref{subsec:AnhangDataAssemblySet}}             
		\\ \hline

		\multicolumn{1}{|L{3.6cm}|}{\cellcolor[HTML]{FFFFFF} ModuleTypePackage: ProcessValueSet.Com-plexTypes V2.0.0 }         & 
		\multicolumn{1}{L{2.8cm}|}{\cellcolor[HTML]{FFFFFF} SUC Struct\-Process\-Value\-In}              &      
		\multicolumn{1}{L{5.43cm}|}{\cellcolor[HTML]{FFFFFF} Schnittstelle zum Einlesen eines Werts mit strukturiertem Datentyp von einer anderen LEA }              &     
		\multicolumn{1}{L{2.45cm}|}{\cellcolor[HTML]{FFFFFF} Anhang~\ref{subsec:AnhangProcessValueSet}}             
		\\ \hline

		\multicolumn{1}{|L{3.6cm}|}{\cellcolor[HTML]{FFFFFF} ModuleTypePackage: ProcessValueSet.Com-plexTypes V2.0.0 }         & 
		\multicolumn{1}{L{2.8cm}|}{\cellcolor[HTML]{FFFFFF} SUC Array\-Process\-Value\-In}              &      
		\multicolumn{1}{L{5.43cm}|}{\cellcolor[HTML]{FFFFFF} Schnittstelle zum Einlesen von Werten eines in einer anderen LEA verwalteten Arrays}              &     
		\multicolumn{1}{L{2.45cm}|}{\cellcolor[HTML]{FFFFFF} Anhang~\ref{subsec:AnhangProcessValueSet}}             
		\\ \hline

		\multicolumn{1}{|L{3.6cm}|}{\cellcolor[HTML]{FFFFFF} ModuleTypePackage: ProcessValueSet.Com-plexTypes V2.0.0 }         & 
		\multicolumn{1}{L{2.8cm}|}{\cellcolor[HTML]{FFFFFF} SUC Output\-Element}              &      
		\multicolumn{1}{L{5.43cm}|}{\cellcolor[HTML]{FFFFFF} Abstrakte Schnittstelle, von der spezifische Prozesswertausgänge verschiedener Datentypen abgeleitet werden können}              &     
		\multicolumn{1}{L{2.45cm}|}{\cellcolor[HTML]{FFFFFF} Anhang~\ref{subsec:AnhangProcessValueSet}}             
		\\ \hline

		\multicolumn{1}{|L{3.6cm}|}{\cellcolor[HTML]{FFFFFF} ModuleTypePackage: ProcessValueSet.Com-plexTypes V2.0.0 }         & 
		\multicolumn{1}{L{2.8cm}|}{\cellcolor[HTML]{FFFFFF} SUC Array\-Process\-Value\-Out}              &      
		\multicolumn{1}{L{5.43cm}|}{\cellcolor[HTML]{FFFFFF} Schnittstelle zur Bereitstellung der Werte eines LEA-intern verwalteten Arrays für eine andere LEA}              &     
		\multicolumn{1}{L{2.45cm}|}{\cellcolor[HTML]{FFFFFF} Anhang~\ref{subsec:AnhangProcessValueSet}}             
		\\ \hline

		\multicolumn{1}{|L{3.6cm}|}{\cellcolor[HTML]{FFFFFF} ModuleTypePackage: DataAssemblySet.\newline ComplexTypes V2.0.0 }         & 
		\multicolumn{1}{L{2.8cm}|}{\cellcolor[HTML]{FFFFFF} SUC Struct\-Man}              &      
		\multicolumn{1}{L{5.43cm}|}{\cellcolor[HTML]{FFFFFF} Schnittstelle zur Manipulation von Werten strukturierter Datentypen durch einen Operator}              &     
		\multicolumn{1}{L{2.45cm}|}{\cellcolor[HTML]{FFFFFF} Anhang~\ref{subsec:AnhangDataAssemblySet}}             
		\\ \hline

		\multicolumn{1}{|L{3.6cm}|}{\cellcolor[HTML]{FFFFFF} ModuleTypePackage: DataAssemblySet.\newline ComplexTypes V2.0.0 }         & 
		\multicolumn{1}{L{2.8cm}|}{\cellcolor[HTML]{FFFFFF} SUC Struct\-Man\-Int}              &      
		\multicolumn{1}{L{5.43cm}|}{\cellcolor[HTML]{FFFFFF} Schnittstelle zur Manipulation von Werten strukturierter Datentypen durch einen Operator oder die LEA selbst}              &     
		\multicolumn{1}{L{2.45cm}|}{\cellcolor[HTML]{FFFFFF} Anhang~\ref{subsec:AnhangDataAssemblySet}}             
		\\ \hline

		\multicolumn{1}{|L{3.6cm}|}{\cellcolor[HTML]{FFFFFF} ModuleTypePackage: DataAssemblySet.\newline ComplexTypes V2.0.0 }         & 
		\multicolumn{1}{L{2.8cm}|}{\cellcolor[HTML]{FFFFFF} SUC Array\-Man}              &      
		\multicolumn{1}{L{5.43cm}|}{\cellcolor[HTML]{FFFFFF} Schnittstelle zur Manipulation von Werten mit Array-Datentypen durch einen Operator}              &     
		\multicolumn{1}{L{2.45cm}|}{\cellcolor[HTML]{FFFFFF} Anhang~\ref{subsec:AnhangDataAssemblySet}}             
		\\ \hline

		\multicolumn{1}{|L{3.6cm}|}{\cellcolor[HTML]{FFFFFF} ModuleTypePackage: DataAssemblySet.\newline ComplexTypes V2.0.0 }         & 
		\multicolumn{1}{L{2.8cm}|}{\cellcolor[HTML]{FFFFFF} SUC Array\-Man\-Int}              &      
		\multicolumn{1}{L{5.43cm}|}{\cellcolor[HTML]{FFFFFF} Schnittstelle zur Manipulation von Werten mit Array-Datentypen durch einen Operator oder die LEA selbst}              &     
		\multicolumn{1}{L{2.45cm}|}{\cellcolor[HTML]{FFFFFF} Anhang~\ref{subsec:AnhangDataAssemblySet}}             
		\\ \hline


		% Modelldefinitionen
		\multicolumn{4}{|L{14.8cm}|}{\cellcolor[HTML]{E0E0E0}\textbf{Modelldefinitionen}} 
		\\ \hline

		\multicolumn{1}{|L{3.6cm}|}{\cellcolor[HTML]{E0E0E0}\textbf{Profil}}        	& 
		\multicolumn{1}{L{2.8cm}|}{\cellcolor[HTML]{E0E0E0}\textbf{Definition}} 			& 
		\multicolumn{1}{L{5.43cm}|}{\cellcolor[HTML]{E0E0E0}\textbf{Beschreibung}} 			& 
		\multicolumn{1}{L{2.45cm}|}{\cellcolor[HTML]{E0E0E0}\textbf{Detail-beschreibung}} 
		\\ \hline

		\multicolumn{1}{|L{3.6cm}|}{\cellcolor[HTML]{FFFFFF} ModuleTypePackage: ServiceSet.Base V2.0.0 }                    & 
		\multicolumn{1}{L{2.8cm}|}{\cellcolor[HTML]{FFFFFF} SUC Service\-Parameter\newline (Erweiterung)}              &      
		\multicolumn{1}{L{5.43cm}|}{\cellcolor[HTML]{FFFFFF} Erweiterung des bestehenden ServiceParameter-Modells um Function\-Classification\-Attributes}              &     
		\multicolumn{1}{L{2.45cm}|}{\cellcolor[HTML]{FFFFFF} Anhang~\ref{subsec:AnhangServiceSet}}             
		\\ \hline

		\multicolumn{1}{|L{3.6cm}|}{\cellcolor[HTML]{FFFFFF} ModuleTypePackage: ServiceSet.Logistics V2.0.0 }         & 
		\multicolumn{1}{L{2.8cm}|}{\cellcolor[HTML]{FFFFFF} SUC Logistics\-Interaction}              &      
		\multicolumn{1}{L{5.43cm}|}{\cellcolor[HTML]{FFFFFF} Modell zur Repräsentation aller logistikspezifischen Anfragen einer LEA an einen LOL}              &     
		\multicolumn{1}{L{2.45cm}|}{\cellcolor[HTML]{FFFFFF} Anhang~\ref{subsec:AnhangServiceSet}}             
		\\ \hline

		\multicolumn{1}{|L{3.6cm}|}{\cellcolor[HTML]{FFFFFF} ModuleTypePackage: ServiceSet.Logistics V2.0.0 }         & 
		\multicolumn{1}{L{2.8cm}|}{\cellcolor[HTML]{FFFFFF} SUC Logistics\-Question}              &      
		\multicolumn{1}{L{5.43cm}|}{\cellcolor[HTML]{FFFFFF} Modell einer LEA-Anfrage an einen LOL}              &     
		\multicolumn{1}{L{2.45cm}|}{\cellcolor[HTML]{FFFFFF} Anhang~\ref{subsec:AnhangServiceSet}}             
		\\ \hline

		\multicolumn{1}{|L{3.6cm}|}{\cellcolor[HTML]{FFFFFF} ModuleTypePackage: ServiceSet.Logistics V2.0.0 }         & 
		\multicolumn{1}{L{2.8cm}|}{\cellcolor[HTML]{FFFFFF} SUC Product\-Parameter\-Request}              &      
		\multicolumn{1}{L{5.43cm}|}{\cellcolor[HTML]{FFFFFF} Modell einer LEA-Anfrage nach produktspezifischen Parametern}              &     
		\multicolumn{1}{L{2.45cm}|}{\cellcolor[HTML]{FFFFFF} Anhang~\ref{subsec:AnhangServiceSet}}             
		\\ \hline

		\multicolumn{1}{|L{3.6cm}|}{\cellcolor[HTML]{FFFFFF} ModuleTypePackage: ServiceSet.Logistics V2.0.0 }         & 
		\multicolumn{1}{L{2.8cm}|}{\cellcolor[HTML]{FFFFFF} SUC Packaging\-Parameter\-Request}              &      
		\multicolumn{1}{L{5.43cm}|}{\cellcolor[HTML]{FFFFFF} Modell einer LEA-Anfrage nach verpackungsspezifischen Parametern}              &     
		\multicolumn{1}{L{2.45cm}|}{\cellcolor[HTML]{FFFFFF} Anhang~\ref{subsec:AnhangServiceSet}}             
		\\ \hline

		\multicolumn{1}{|L{3.6cm}|}{\cellcolor[HTML]{FFFFFF} ModuleTypePackage: ServiceSet.Logistics V2.0.0 }         & 
		\multicolumn{1}{L{2.8cm}|}{\cellcolor[HTML]{FFFFFF} SUC Product\-Parameter\-Updated\-Info}              &      
		\multicolumn{1}{L{5.43cm}|}{\cellcolor[HTML]{FFFFFF} Modell einer LEA-Anfrage zur Information des LOL nach der Änderung eines produktspezifischen Parametersatzes }              &     
		\multicolumn{1}{L{2.45cm}|}{\cellcolor[HTML]{FFFFFF} Anhang~\ref{subsec:AnhangServiceSet}}             
		\\ \hline

		\multicolumn{1}{|L{3.6cm}|}{\cellcolor[HTML]{FFFFFF} ModuleTypePackage: ServiceSet.Logistics V2.0.0 }         & 
		\multicolumn{1}{L{2.8cm}|}{\cellcolor[HTML]{FFFFFF} SUC Packaging\-Parameter\-Updated\-Info}              &      
		\multicolumn{1}{L{5.43cm}|}{\cellcolor[HTML]{FFFFFF} Modell einer LEA-Anfrage zur Information des LOL nach der Änderung eines verpackungsspezifischen Parametersatzes }              &     
		\multicolumn{1}{L{2.45cm}|}{\cellcolor[HTML]{FFFFFF} Anhang~\ref{subsec:AnhangServiceSet}}             
		\\ \hline

		\multicolumn{1}{|L{3.6cm}|}{\cellcolor[HTML]{FFFFFF} ModuleTypePackage: ServiceSet.Logistics V2.0.0 }                    & 
		\multicolumn{1}{L{2.8cm}|}{\cellcolor[HTML]{FFFFFF} RC Has\-Logistics\-Interaction}              &      
		\multicolumn{1}{L{5.43cm}|}{\cellcolor[HTML]{FFFFFF} Modell zur Zuordnung einer LogisticsInteraction zum Dienst einer LEA}              &     
		\multicolumn{1}{L{2.45cm}|}{\cellcolor[HTML]{FFFFFF} Anhang~\ref{subsec:AnhangServiceSet}}             
		\\ \hline


		\multicolumn{1}{|L{3.6cm}|}{\cellcolor[HTML]{FFFFFF} ModuleTypePackage: ServiceSet.Logistics V2.0.0 }                    & 
		\multicolumn{1}{L{2.8cm}|}{\cellcolor[HTML]{FFFFFF} SUC Service\newline (Erweiterung)}              &      
		\multicolumn{1}{L{5.43cm}|}{\cellcolor[HTML]{FFFFFF} Erweiterung des bestehenden Service-Modells um die Möglichkeit, die RC HasLogisticsInteraction anzufügen}              &     
		\multicolumn{1}{L{2.45cm}|}{\cellcolor[HTML]{FFFFFF} Anhang~\ref{subsec:AnhangServiceSet}}             
		\\ \hline

	\end{longtable}
\end{footnotesize}

Über diese Modell- und Schnittstellendefinitionen hinaus werden \emph{FunctionClassificationAttributes} für CES- und SES-Prozeduren sowie für die Parameter \emph{ProductId}, \emph{LogisticsObjectStatus}, \emph{ProductDataSet}, \emph{PackagingId} und \emph{PackagingDataSet} eingeführt (Anhang~\ref{subsec:AnhangServiceSet}). Zudem wird in Anhang~\ref{subsec:AnhangServerAssemblySet} ein Mechanismus zur Abbildung komplexer Datentypen in einem OPC~UA-Server beschrieben. Weiterhin zeigte sich, dass die bestehenden MTP-Schnittstellendefinitionen für die Automatisierung von LEAs teils zu umfangreich und komplex sind. Eine Möglichkeit zur Komplexitätsreduktion ist in Anhang~\ref{subsec:Komplexitaetsreduktion} zu finden. 


